\chapter{Statistical Foundations}
\label{ch:stats}

\begin{quote}\itshape
``The theory of probabilities is at bottom nothing but common sense reduced to
calculus.'' --- P.-S.~Laplace.
\end{quote}

\noindent
The chapters that follow lean on statistics at every turn: a response surface is
\emph{fitted}, an effect is judged \emph{significant}, a design is called
\emph{efficient}, a fermentation is described as \emph{variable} but a probe as
\emph{uncertain}. Each of these words has a precise technical meaning, and a
reader fluent in them will find the later mathematics not merely legible but
inevitable. This chapter assembles that vocabulary from the ground up. It is
deliberately self-contained---a working physical scientist who has not opened a
statistics text since university should be able to read it straight
through---but it is also curated: we develop \emph{exactly} the ideas that
Chapters~\ref{ch:design} and~\ref{ch:ferment} use, and no more.

Three threads run through the chapter. The first is the \emph{inferential leap}:
we never see the truth, only a finite, noisy sample of it, and statistics is the
disciplined art of saying what the sample does and does not tell us about the
truth. The second is that every method here is, at bottom, an \emph{algorithm}
one can write in a dozen lines; wherever a formula appears we also give the
recipe that computes it, and point to the small teaching module
(\code{downstream\_doe.foundations.stats\_demo}) from which this chapter's
figures are in fact generated. The third is the distinction---fundamental to the
whole monograph---between \emph{variability}, the irreducible spread of a real
process, and \emph{uncertainty}, the fog of imperfect measurement.

%==============================================================================
\section{Populations, samples, and the inferential leap}
\label{sec:stats-pop}

Imagine every yogurt batch that a process \emph{could ever} produce under fixed
settings. That infinite, hypothetical collection is the \emph{population}; the
number we care about---say its true mean set-time---is a fixed but unknown
\emph{parameter}, written with a Greek letter such as $\mu$. We cannot run
infinitely many batches. We run a handful: a \emph{sample}
$x_1,\dots,x_n$. Any number we compute from the sample, such as the sample mean
$\bar x = \tfrac1n\sum_i x_i$, is a \emph{statistic}, and a statistic used to
guess a parameter is an \emph{estimator}, written with a hat: $\hat\mu = \bar x$.

The central fact---the inferential leap---is that the estimator is itself
\emph{random}. Run a different five batches and $\bar x$ comes out different. An
estimator therefore has its own distribution, the \emph{sampling distribution},
and the whole of classical statistics is the study of how far an estimate is
liable to sit from the parameter it estimates. An estimator is \emph{unbiased}
if it is right \emph{on average}, $\E[\hat\mu]=\mu$; its scatter about that
average is what a \emph{standard error} (Section~\ref{sec:stats-est}) measures.

\begin{keybox}{The four words to keep straight}
\emph{Population/parameter} (the unknowable truth, $\mu$) versus
\emph{sample/statistic} (what we compute, $\bar x$). A \emph{parameter} is fixed;
a \emph{statistic} is random because the sample is random. Estimation is the
bridge from the second to the first, and it is never free of error---it can only
be quantified.
\end{keybox}

%==============================================================================
\section{The language of distributions}
\label{sec:stats-dist}

A \emph{random variable} $X$ is a quantity whose value is subject to chance.
A \emph{continuous} random variable is described by its \emph{probability
density function} (PDF) $f(x)$, which is non-negative and integrates to one; the
probability of landing in an interval is the area under the density,
$P(a\le X\le b)=\int_a^b f(x)\dd x$. Its running integral is the
\emph{cumulative distribution function} (CDF) $F(x)=P(X\le x)=\int_{-\infty}^x
f$, which climbs monotonically from $0$ to $1$ and answers ``how probable is a
value no larger than $x$?'' (Figure~\ref{fig:stats-dist}).

Two numbers summarise a distribution. The \emph{mean} (or expectation) is its
balance point,
%
\begin{equation}
\mu = \E[X] = \int x\,f(x)\dd x,
\end{equation}
%
and the \emph{variance} is the mean squared distance from that balance point,
%
\begin{equation}
\sigma^2 = \Var(X) = \E\!\bigl[(X-\mu)^2\bigr]
        = \E[X^2]-\mu^2,
\label{eq:variance}
\end{equation}
%
whose square root $\sigma$, the \emph{standard deviation}, has the same units as
$X$ and is the natural ``width'' of the distribution. When a spread is more
meaningful relative to the mean---as for a strictly positive quantity like a rate
constant---one reports the dimensionless \emph{coefficient of variation}
$\mathrm{CV}=\sigma/\mu$.

\paragraph{The normal distribution.}
The distribution that recurs everywhere is the \emph{normal} (or Gaussian),
%
\begin{equation}
f(x) = \frac{1}{\sigma\sqrt{2\pi}}
       \exp\!\left(-\frac{(x-\mu)^2}{2\sigma^2}\right),
\label{eq:normal}
\end{equation}
%
the symmetric bell of Figure~\ref{fig:stats-dist}(a), with about $68\%$ of its
mass within $\pm\sigma$ of the mean and $95\%$ within $\pm2\sigma$. Its
ubiquity is not an accident but a theorem: the \emph{central limit theorem}
states that a sum (or average) of many small, independent contributions is
approximately normal whatever the shape of the individual pieces. This is why
measurement noise---the accumulation of countless tiny perturbations---is so
often modelled as Gaussian, as it is for every virtual instrument in this
package.

\paragraph{The lognormal distribution.}
When a quantity must stay positive and varies \emph{multiplicatively}---each
batch a random \emph{factor} times a nominal value, rather than a random
\emph{offset}---the right model is the \emph{lognormal}: $X$ is lognormal if
$\ln X$ is normal. This is precisely the law Chapter~\ref{ch:ferment} uses to
jitter physical parameters between batches, $\theta = \bar\theta\,
e^{\sigma_\theta Z}$ with $Z\sim\mathcal N(0,1)$, because a growth rate or a
buffering capacity can be $20\%$ high or low but never negative, and for small
$\sigma_\theta$ the parameter $\sigma_\theta$ is itself the coefficient of
variation.

\begin{figure}[t]
\centering
\includegraphics[width=\linewidth]{foundations_distributions.png}
\caption[Density, distribution, and a finite sample]{The anatomy of a
continuous random variable, here a normal with $\mu=4.6$, $\sigma=0.4$.
\emph{(a)} The density $f(x)$: probability is area, and the shaded bands hold
$68\%$ ($\pm\sigma$) and $95\%$ ($\pm2\sigma$) of it. \emph{(b)} The CDF
$F(x)=P(X\le x)$, from which any interval probability is read directly.
\emph{(c)} A finite sample of $400$ draws: the histogram approximates the true
density and the sample mean $\bar x$ approximates $\mu$---the closer, the larger
the sample. Generated by \code{stats\_demo} via the package's RNG.}
\label{fig:stats-dist}
\end{figure}

%==============================================================================
\section{Estimation and the standard error}
\label{sec:stats-est}

Given a sample, the natural estimators are the \emph{sample mean} and
\emph{sample variance}
%
\begin{equation}
\bar x = \frac1n\sum_{i=1}^n x_i,
\qquad
s^2 = \frac{1}{n-1}\sum_{i=1}^n (x_i-\bar x)^2 .
\label{eq:sample-stats}
\end{equation}
%
The divisor $n-1$ rather than $n$---\emph{Bessel's correction}---makes $s^2$
unbiased; intuitively, one degree of freedom is spent estimating $\bar x$, so
only $n-1$ independent residuals remain. \emph{Degrees of freedom} (the number
of independent pieces of information left after fitting) is a bookkeeping
quantity we shall meet again in regression and the $F$-test.

How precise is $\bar x$? Because the sample mean averages $n$ independent draws,
its variance is $\sigma^2/n$, so its standard deviation---called the
\emph{standard error} to distinguish it from the spread of the raw data---is
%
\begin{equation}
\mathrm{SE}(\bar x) = \frac{\sigma}{\sqrt n}\;\approx\;\frac{s}{\sqrt n}.
\label{eq:se}
\end{equation}
%
The $\sqrt n$ is the central, and sobering, economics of experimentation: to
halve the error of a mean one must \emph{quadruple} the number of runs. From the
standard error follows the \emph{confidence interval}, the interval that, over
many hypothetical repetitions of the whole experiment, would contain the true
$\mu$ a stated fraction (say $95\%$) of the time:
%
\begin{equation}
\bar x \;\pm\; t_{0.975,\,n-1}\,\frac{s}{\sqrt n},
\label{eq:ci}
\end{equation}
%
with $t$ the critical value of Student's $t$-distribution (which widens the
interval to pay for using $s$ in place of the unknown $\sigma$, and tends to the
normal value $1.96$ as $n$ grows). The subtle, often-misread meaning is worth
stating plainly: the $95\%$ attaches to the \emph{procedure}, not to a single
interval; the truth is fixed and it is the interval that is random.

%==============================================================================
\section{Fitting a model by least squares}
\label{sec:stats-ols}

Most of Chapter~\ref{ch:design} is about relating an output $y$ to inputs
$x_1,\dots,x_p$ through a \emph{linear model}
%
\begin{equation}
y = \beta_0 + \beta_1 x_1 + \dots + \beta_p x_p + \eps,
\qquad \eps\sim\mathcal N(0,\sigma^2),
\label{eq:linmodel}
\end{equation}
%
where the \emph{coefficients} $\beta_j$ are unknown and $\eps$ is random error.
``Linear'' refers to linearity in the \emph{coefficients}, not the inputs: an
interaction $x_1x_2$ or a curvature term $x_1^2$ is just another column, which is
why the same machinery fits a full quadratic response surface. Stacking the $n$
observations, the model is $\bm y = \bm X\bm\beta + \bm\eps$, where the
\emph{design matrix} $\bm X$ has one row per run and one column per term.

\emph{Least squares} chooses the coefficients that minimise the sum of squared
\emph{residuals} $r_i = y_i-\hat y_i$, the vertical gaps between data and fit
(Figure~\ref{fig:stats-ols}). Setting the gradient of $\sum_i r_i^2$ to zero
gives the \emph{normal equations}, whose solution and covariance are the
classical
%
\begin{equation}
\hat{\bm\beta} = (\bm X\transp\bm X)^{-1}\bm X\transp\bm y,
\qquad
\Var(\hat{\bm\beta}) = \sigma^2(\bm X\transp\bm X)^{-1}.
\label{eq:normaleq}
\end{equation}
%
The diagonal of that covariance gives the \emph{standard error} of each
coefficient, and hence a confidence interval and significance test for every
effect. Two properties of $\bm X\transp\bm X$ matter enormously in design. If the
columns are \emph{orthogonal} (uncorrelated), $\bm X\transp\bm X$ is diagonal,
every coefficient is estimated independently and with minimum variance, and
dropping a term does not change the others---this is exactly why the factorial
designs of Section~\ref{sec:factorial} code their factors to $\pm1$. If instead
the columns are nearly \emph{collinear}, $\bm X\transp\bm X$ is close to singular,
its inverse explodes, and the coefficients become wildly uncertain---the
pathology that motivates the latent-variable methods of Section~\ref{sec:mv}.

The quality of the fit is summarised by the \emph{coefficient of determination}
%
\begin{equation}
R^2 = 1 - \frac{\sum_i (y_i-\hat y_i)^2}{\sum_i (y_i-\bar y)^2},
\label{eq:r2}
\end{equation}
%
the fraction of the response's variance the model explains: $R^2=1$ is a perfect
fit, $R^2=0$ no better than the mean.

\begin{keybox}{Algorithm: least squares by the normal equations}
The whole of Eq.~\eqref{eq:normaleq} is four lines of code
(\code{stats\_demo.ols\_fit}):
\begin{verbatim}
    XtX  = X.T @ X                      # the normal matrix
    beta = np.linalg.solve(XtX, X.T@y)  # solve, don't invert
    resid   = y - X @ beta
    sigma2  = resid @ resid / (n - p)   # residual variance
    cov     = sigma2 * np.linalg.inv(XtX)
\end{verbatim}
One solves the system rather than forming $(\bm X\transp\bm X)^{-1}$ explicitly
(it is both faster and numerically better); the inverse is computed only to
report the covariance, from which every standard error follows.
\end{keybox}

\begin{figure}[t]
\centering
\includegraphics[width=0.75\linewidth]{foundations_ols.png}
\caption[Least squares and its confidence band]{Ordinary least squares fits the
line that minimises the total squared length of the grey residual stems. The
shaded band is the $95\%$ confidence interval for the \emph{mean} response,
computed from the coefficient covariance of Eq.~\eqref{eq:normaleq}; it is
narrowest at the centre of the data and flares outward, the geometric signature
of extrapolation risk. Generated by \code{stats\_demo.ols\_fit}.}
\label{fig:stats-ols}
\end{figure}

%==============================================================================
\section{Generalized linear models}
\label{sec:stats-glm}

Ordinary least squares quietly assumes that the response is a continuous
quantity that may range over the whole real line, with normal, constant-variance
scatter. Much experimental data is not like that. The outcome of a screening run
may be a \emph{yes/no}---did the blend reach the set point? did the pool meet the
purity specification?---or a \emph{count}---how many impurity peaks appeared?---or
a \emph{proportion} bounded in $[0,1]$. Fitting a straight line to a $0/1$
outcome is not merely inelegant; it predicts probabilities below zero and above
one (Figure~\ref{fig:stats-glm}), and its constant-variance assumption is simply
false for such data.

The \emph{generalized linear model} (GLM)~\cite{mccullagh} keeps everything that
is good about the linear model and repairs exactly what fails. It retains the
\emph{linear predictor} $\eta = \bm x\transp\bm\beta$---the weighted sum of
inputs---but inserts a \emph{link function} $g$ between that predictor and the
mean of the response,
%
\begin{equation}
g(\mu) = \eta = \bm x\transp\bm\beta,
\qquad
\mu = \E[y\mid\bm x],
\label{eq:glm}
\end{equation}
%
and lets the response follow whatever distribution suits its nature (Bernoulli
for yes/no, Poisson for counts, normal for the classical case, which is just the
GLM with an identity link). For a binary outcome the natural choice is the
\emph{logit} link, giving \emph{logistic regression}: the log-odds are linear in
the inputs, and inverting the link returns a probability guaranteed to lie in
$[0,1]$,
%
\begin{equation}
\log\frac{\mu}{1-\mu} = \bm x\transp\bm\beta
\quad\Longleftrightarrow\quad
\mu = \sigma(\bm x\transp\bm\beta)
    = \frac{1}{1+e^{-\bm x\transp\bm\beta}} .
\label{eq:logistic}
\end{equation}
%
A coefficient $\beta_j$ now has a clean reading: it is the change in the
\emph{log-odds} of the event per unit of $x_j$, so $e^{\beta_j}$ is an
\emph{odds ratio}. The same template, with a $\log$ link and a Poisson response,
gives \emph{Poisson regression} for count data.

\paragraph{Fitting: maximum likelihood and the return of least squares.}
There is no closed form like Eq.~\eqref{eq:normaleq}, because the link is
nonlinear. One instead chooses $\bm\beta$ to maximise the \emph{likelihood}---the
probability the model assigns to the observed data---and the elegant result is
that the optimisation reduces to solving a \emph{weighted} least-squares problem
over and over. This is \emph{iteratively reweighted least squares} (IRLS): at
each step the current fit defines a weight $w_i=\mu_i(1-\mu_i)$ for each
observation (largest where the prediction is most uncertain, near
$\mu=\tfrac12$), and the coefficients are updated by the weighted normal
equations.

\begin{keybox}{Algorithm: logistic regression by IRLS}
The whole of GLM fitting is the OLS update of Section~\ref{sec:stats-ols} with
weights, iterated to convergence (\code{stats\_demo.logistic\_fit}):
\begin{verbatim}
    mu   = sigmoid(X @ beta)            # current fitted probabilities
    w    = mu * (1 - mu)                # IRLS weights
    XtWX = X.T @ (w[:, None] * X)       # weighted normal matrix
    beta += np.linalg.solve(XtWX, X.T @ (y - mu))   # Newton step
\end{verbatim}
Swap the link and the weight rule and the same loop fits Poisson, probit, and
the rest of the GLM family. The in-sample fit is judged not by $R^2$ but by the
\emph{deviance}, the GLM's analogue of the residual sum of squares.
\end{keybox}

GLMs are the workhorse for \emph{analysing} experimental data whose response is
not a tidy Gaussian. In Chapter~\ref{ch:ferment} we use a (regularized) logistic
regression to ask which strains make a blend likely to set
(Section~\ref{sec:ferm-analysis}). The same idea reaches naturally into
downstream chromatography---a logistic model of whether a run meets its purity
specification, or a Poisson model of the number of impurity peaks in a
chromatogram---which Section~\ref{sec:glm} develops into a \emph{probabilistic
design space}.

\begin{figure}[t]
\centering
\includegraphics[width=0.72\linewidth]{foundations_glm.png}
\caption[A generalized linear model]{Why a binary outcome needs a GLM. The
points are $0/1$ responses (did the event occur?) against a predictor. A linear
fit (red dashed) sails out of the shaded forbidden zones, predicting
probabilities below $0$ and above $1$; the logistic GLM (blue) passes the linear
predictor through the logit link, so its fitted probability is always a genuine
probability. Generated by \code{stats\_demo.logistic\_fit}.}
\label{fig:stats-glm}
\end{figure}

%==============================================================================
\section{Is the effect real? Hypothesis tests and the analysis of variance}
\label{sec:stats-anova}

Having estimated an effect, we must ask whether it could plausibly be an
artefact of noise. The classical apparatus is the \emph{hypothesis test}. One
states a \emph{null hypothesis} $H_0$---``this effect is zero''---computes a
\emph{test statistic} whose distribution under $H_0$ is known, and reports the
\emph{$p$-value}: the probability, \emph{if $H_0$ were true}, of seeing a
statistic at least as extreme as the one observed. A small $p$-value
(conventionally $p<0.05$) means the data are hard to reconcile with ``no
effect,'' so we reject $H_0$. The $p$-value is \emph{not} the probability that
$H_0$ is true---a perennial confusion---but a statement about the data under an
assumption.

The \emph{analysis of variance} (ANOVA) is the form this takes for the linear
model, and the engine of Section~\ref{sec:factorial}. Its insight is that the
total scatter of the response decomposes exactly into a part the model explains
and an unexplained residual,
%
\begin{equation}
\underbrace{\sum_i(y_i-\bar y)^2}_{\text{SS}_{\text{total}}}
=
\underbrace{\sum_i(\hat y_i-\bar y)^2}_{\text{SS}_{\text{model}}}
+
\underbrace{\sum_i(y_i-\hat y_i)^2}_{\text{SS}_{\text{residual}}} .
\label{eq:ssdecomp}
\end{equation}
%
Dividing each \emph{sum of squares} by its degrees of freedom gives a \emph{mean
square}, and the ratio of the two,
%
\begin{equation}
F = \frac{\text{SS}_{\text{model}}/\mathrm{df}_{\text{model}}}
         {\text{SS}_{\text{residual}}/\mathrm{df}_{\text{residual}}},
\label{eq:ftest}
\end{equation}
%
is, under the null hypothesis that the model adds nothing, a draw from the
$F$-distribution. A large $F$ (small $p$) says the explained variation is too
big to be chance. Applied term by term, this is exactly how the package decides
which factors and interactions matter (\code{stats\_demo.regression\_anova}, and
at scale the \code{anova\_lm} of Section~\ref{sec:factorial}); the same
sum-of-squares partition reappears, in Section~\ref{sec:stats-varcomp}, as the
way to split variability from uncertainty.

%==============================================================================
\section{Honest prediction: overfitting and cross-validation}
\label{sec:stats-cv}

There is a trap in $R^2$: it can always be driven toward one by adding terms,
because a flexible enough model can memorise the very noise it was fitted to.
Such a model \emph{overfits}---it describes the sample it has seen and predicts
new data badly. The cure is to judge a model on data it did \emph{not} see.

\emph{$K$-fold cross-validation} does this with the data in hand. The runs are
split into $K$ equal folds; each fold is held out in turn, the model is fitted
on the other $K-1$, and the held-out fold is predicted. The
cross-validated coefficient,
%
\begin{equation}
R^2_{\mathrm{CV}} = 1 - \frac{\sum_i\bigl(y_i-\hat y_{-k(i)}(\bm x_i)\bigr)^2}
                              {\sum_i (y_i-\bar y)^2},
\label{eq:cv}
\end{equation}
%
where $\hat y_{-k(i)}$ is the prediction for run $i$ from a model trained without
its fold, measures genuine \emph{out-of-sample} skill. Unlike the in-sample
$R^2$, it \emph{falls} once added complexity begins to fit noise, giving an
honest stopping rule. This single idea appears under two names in the chapters
that follow---the $Q^2$ of partial least squares (Section~\ref{sec:mv}) and the
$R^2_{\mathrm{CV}}$ that grades the tree ensembles of
Section~\ref{sec:ferm-analysis}---but it is always Eq.~\eqref{eq:cv}.

%==============================================================================
\section{Regularization: taming complexity and collinearity}
\label{sec:stats-reg}

Cross-validation \emph{diagnoses} overfitting; \emph{regularization} prevents it.
The setting is the one that defeats ordinary least squares: more candidate
effects than runs ($p>n$), or columns so collinear that $\bm X\transp\bm X$ is
near-singular and the coefficient variances of Eq.~\eqref{eq:normaleq} explode.
Both are everyday in screening, where one may wish to entertain hundreds of
strains or interaction terms from a couple of hundred experiments. The remedy is
to stop minimising the residual sum of squares alone and instead minimise it
\emph{plus a penalty on the size of the coefficients},
%
\begin{equation}
\hat{\bm\beta} = \arg\min_{\bm\beta}\;
   \underbrace{\sum_i\bigl(y_i-\bm x_i\transp\bm\beta\bigr)^2}_{\text{fit}}
   \;+\; \lambda\,\underbrace{P(\bm\beta)}_{\text{penalty}},
\label{eq:penalised}
\end{equation}
%
where $\lambda\ge0$ sets how hard the penalty bites. The penalty buys a little
\emph{bias} in exchange for a large cut in \emph{variance}---the
\emph{bias--variance trade-off}---and the net effect on out-of-sample error is
usually a clear win. The two classic penalties differ in their geometry:

\begin{itemize}
  \item \textbf{Ridge} ($L_2$): $P(\bm\beta)=\sum_j\beta_j^2$. This adds
  $\lambda$ to the diagonal of the normal matrix,
  $\hat{\bm\beta}=(\bm X\transp\bm X+\lambda\bm I)^{-1}\bm X\transp\bm y$, which
  is \emph{always} invertible---curing collinearity outright. Ridge
  \emph{shrinks} every coefficient smoothly toward zero but sets none exactly to
  zero (Figure~\ref{fig:stats-reg}, left).
  \item \textbf{Lasso}~\cite{tibshirani} ($L_1$): $P(\bm\beta)=\sum_j|\beta_j|$. The absolute-value
  penalty has a corner at the origin, and that corner does something ridge
  cannot: it drives the coefficients of unhelpful inputs \emph{exactly} to zero
  (Figure~\ref{fig:stats-reg}, right). The lasso therefore performs
  \emph{variable selection}---its fitted model names a sparse, interpretable
  subset of the inputs that matter. The \emph{elastic net} mixes the two,
  inheriting the lasso's sparsity and the ridge's grace under correlated columns.
\end{itemize}

The penalty strength $\lambda$ is the one knob, and it is chosen by exactly the
cross-validation of Section~\ref{sec:stats-cv}: sweep $\lambda$, and keep the
value with the best out-of-sample score.

\begin{keybox}{Algorithm: ridge in one line, the lasso by coordinate descent}
Ridge is a one-line modification of the normal equations
(\code{stats\_demo.ridge\_fit}):
\begin{verbatim}
    beta = np.linalg.solve(X.T@X + lam*np.eye(p), X.T@y)
\end{verbatim}
The lasso has no closed form, but cycling over coefficients and
\emph{soft-thresholding} each---shrinking it toward zero and clamping to exactly
zero once the penalty outweighs its pull---converges quickly
(\code{stats\_demo.lasso\_fit}):
\begin{verbatim}
    rho     = X[:, j] @ (y - X@beta + X[:, j]*beta[j])   # partial fit
    beta[j] = sign(rho) * max(abs(rho) - lam, 0) / ||X[:, j]||^2
\end{verbatim}
The $\max(\cdot,0)$ is the corner that creates sparsity.
\end{keybox}

\paragraph{Regularization beyond the linear model.}
The idea is far more general than ridge and the lasso. \emph{Every} method in
this book that fights overfitting is regularizing in disguise. The
\emph{random forest} of Section~\ref{sec:ferm-analysis} regularizes
\emph{implicitly}: by averaging many trees, each grown on a bootstrap resample
and restricted at each split to a random subset of inputs, it trades a little
bias for a large variance reduction---the same bargain as
Eq.~\eqref{eq:penalised}, struck by randomisation rather than a penalty term.
\emph{Gradient boosting} (XGBoost) regularizes \emph{explicitly}: its objective
carries a penalty on the number of leaves and the size of the leaf weights
(the $\Omega(f)$ term of Eq.~\eqref{eq:xgb}), which is Eq.~\eqref{eq:penalised}
transplanted onto trees. Seeing regularization as one idea---pay a little bias to
buy a lot of variance---unifies the linear and the ensemble methods that
Chapter~\ref{ch:ferment} sets side by side~\cite{esl}.

\begin{figure}[t]
\centering
\includegraphics[width=\linewidth]{foundations_regularization.png}
\caption[Ridge and lasso regularization paths]{Coefficient paths as the penalty
$\lambda$ relaxes from strong (left) to weak (right); three inputs truly matter
(red) and five are noise (grey). \emph{Left:} ridge shrinks every coefficient
smoothly and in concert, but no coefficient is ever exactly zero. \emph{Right:}
the lasso holds the noise coefficients at \emph{exactly} zero until the penalty
is quite weak, selecting the three real effects---variable selection for free.
Generated by \code{stats\_demo.coefficient\_path}.}
\label{fig:stats-reg}
\end{figure}

%==============================================================================
\section{Error bars without a formula: the bootstrap}
\label{sec:stats-boot}

The standard error~\eqref{eq:se} has a tidy formula because the mean is simple
and the sample large. For a complicated statistic, or a small, skewed sample, no
such formula may exist. The \emph{bootstrap} of Efron supplies one by brute
force, on a single, almost impudent idea: \emph{the sample is the best estimate
of the population we have, so resample from it}.

Concretely, to put an error bar on any statistic $T$ one draws many new samples
of size $n$ \emph{from the observed data, with replacement}, recomputes $T$ on
each, and reads the spread of those replicates as the sampling distribution of
$T$. A confidence interval is then just the empirical quantiles of the
replicates (Figure~\ref{fig:stats-boot}).

\begin{keybox}{Algorithm: the percentile bootstrap}
For \code{stats\_demo.bootstrap\_ci(sample, statistic, n\_boot, alpha)}:
\begin{enumerate}
  \item Repeat \code{n\_boot} times: draw $n$ indices uniformly with replacement
  from $\{1,\dots,n\}$, form the resample, and store $T$ of it.
  \item Report the observed $T$ as the estimate, and the empirical
  $[\alpha/2,\,1-\alpha/2]$ quantiles of the stored replicates as the confidence
  interval.
\end{enumerate}
No distributional assumption, no derivative, no formula for the variance of
$T$---just resampling and counting. This is a special case of the broader
\emph{Monte~Carlo} principle: replace an intractable expectation by an average
over many random draws, the same principle that underlies the uncertainty
propagation of this package.
\end{keybox}

\begin{figure}[t]
\centering
\includegraphics[width=\linewidth]{foundations_bootstrap.png}
\caption[The bootstrap]{The bootstrap turns one sample into a sampling
distribution. \emph{(a)} A small ($n=18$), skewed sample---exactly the case where
a textbook normal interval is unreliable. \emph{(b)} Resampling it $20{,}000$
times with replacement and recomputing the mean yields the sampling distribution
of $\bar x$; its central $95\%$ (green) is the bootstrap confidence interval.
Generated by \code{stats\_demo.bootstrap\_ci}.}
\label{fig:stats-boot}
\end{figure}

%==============================================================================
\section{Variability versus uncertainty}
\label{sec:stats-varcomp}

The most important conceptual distinction in this monograph is between two kinds
of randomness that a single number like ``$\pm0.3$'' hopelessly conflates.
\emph{Variability} (the statisticians' \emph{aleatoric} uncertainty) is real,
physical spread: genuinely identical settings still give different batches,
because the milk lot, the inoculum, and the incubator are never quite the same.
It does not shrink with a better instrument---only with better process control.
\emph{Uncertainty} (\emph{epistemic}) is our imperfect knowledge: measurement
noise, finite samples, a model that is only approximate. It \emph{does} shrink
with more or better data, and can be averaged away by replication.

These can be separated when a design measures several batches, each more than
once. A balanced one-way \emph{variance-components} analysis splits the total
scatter into a \emph{within-batch} part (repeat measurements of the same
batch---this is uncertainty) and a \emph{between-batch} part (the spread of the
batch means above and beyond that noise---this is variability),
%
\begin{equation}
\hat\sigma^2_{\text{within}} = \mathrm{MS}_{\text{within}},
\qquad
\hat\sigma^2_{\text{between}}
   = \frac{\mathrm{MS}_{\text{between}}-\mathrm{MS}_{\text{within}}}{m},
\label{eq:varcomp}
\end{equation}
%
from the same mean squares as the ANOVA, with $m$ the number of replicates per
batch (\code{stats\_demo.variance\_components}). Figure~\ref{fig:stats-var} shows
the picture: tight clusters of replicate measurements (small within-batch noise)
whose \emph{centres} nevertheless scatter widely (large between-batch
variability). Reading a real process this way---and refusing to collapse the two
numbers into one---is the statistical backbone of the ``three faces of
randomness'' that Chapter~\ref{ch:ferment} builds into its stochastic
fermentation model.

\begin{figure}[t]
\centering
\includegraphics[width=0.8\linewidth]{foundations_variance.png}
\caption[Variability versus uncertainty]{Eight batches, each measured six times.
The replicate measurements within a batch cluster tightly---small
\emph{within-batch} noise, which a better probe would shrink---yet the batch
means (red bars) scatter far more widely about the grand mean (dashed)---real
\emph{between-batch} variability, which only better process control can reduce.
The variance-components estimator of Eq.~\eqref{eq:varcomp} recovers the two
separately. Generated by \code{stats\_demo.variance\_components}.}
\label{fig:stats-var}
\end{figure}

%==============================================================================
\section{Why design experiments?}
\label{sec:stats-doe}

Everything so far has taken the data as given and asked what may be inferred from
it. But the scientist usually gets to choose \emph{which} experiments to run, and
that choice---made before a single measurement---governs how much the data can
ever tell. This is the province of the \emph{design of experiments} (DoE), and
its premise, due to Fisher~\cite{fisher}, is bracing: \emph{no analysis, however
clever, can rescue a badly designed experiment.} A statistician consulted afterwards, in
Fisher's famous phrase, can often only conduct a post-mortem. Design is what
turns experimentation from anecdote-gathering into measurement.

\subsection{The error budget of an experiment}

To see what design must control, decompose the error in any measured response
into two fundamentally different kinds.

\emph{Systematic error}, or \emph{bias}, pushes every measurement in the same
direction: a mis-calibrated probe reading half a unit high, a reagent that
degrades over the day so that later runs differ from earlier ones for a reason
that has nothing to do with the factors. Bias does not average out with more
runs---repeating a biased experiment just gives a more precise wrong answer---and
it is the more dangerous error precisely because it hides.

\emph{Random error} scatters measurements unsystematically about the truth, and
it is itself the sum of the two components Section~\ref{sec:stats-varcomp} taught
us to separate: \emph{measurement uncertainty} (the probe noise, the assay
imprecision---epistemic, shrinkable by replication) and \emph{process
variability} (genuine batch-to-batch biology---aleatoric, shrinkable only by
better control). Random error does average out: the standard
error~\eqref{eq:se} falls as $1/\sqrt n$.

\begin{keybox}{The two errors and the tactics that defeat them}
\emph{Systematic error (bias)} $\to$ defeated by \emph{comparison} (run controls
side by side so the bias cancels in the difference) and \emph{randomization}
(below). \emph{Random error} $\to$ defeated by \emph{replication} (averaging
shrinks it as $1/\sqrt n$ and, as a bonus, \emph{measures} it). Confusing the two
is the classic experimental blunder: no number of replicates fixes a bias, and no
calibration fixes irreducible variability.
\end{keybox}

\subsection{Fisher's principles}

Three design principles, all Fisher's, target this error budget directly.

\emph{Replication}---running the same condition more than once---does double duty:
it shrinks the standard error of each estimated effect, and it furnishes a
model-free estimate of the random error against which effects are judged
significant (the residual mean square of the ANOVA,
Section~\ref{sec:stats-anova}). \emph{Randomization}---assigning run conditions to
time slots in random order---is Fisher's quiet masterstroke: it converts an
unknown systematic error (a drift, a warming room) into part of the random error,
so that it inflates the error bars honestly rather than masquerading as a real
effect. \emph{Blocking}---grouping runs into homogeneous blocks (one milk lot, one
day) and comparing \emph{within} them---removes a known nuisance source of
variability from the comparison entirely, sharpening the contrasts of interest.

\subsection{What a good design buys}

Beyond controlling error, the \emph{structure} of a design determines what can be
estimated and how precisely. A \emph{factorial} design---varying several factors
together rather than one-at-a-time---is the decisive idea, for a reason the
one-factor-at-a-time (OFAT) tradition misses entirely: only by varying factors
\emph{jointly} can one detect an \emph{interaction}, where the effect of one
factor depends on the level of another. Figure~\ref{fig:stats-doe} shows the
trap. On a response with an interaction, OFAT---optimise factor one, fix it,
then optimise factor two---climbs to a point and stops, blind to the diagonal
ridge along which the true optimum lies; a factorial, sampling the corners,
sees the tilt at once. Factorial designs are also \emph{efficient} (every run
informs every main effect, so the information per experiment is maximal) and,
when the factors are coded as in Section~\ref{sec:stats-ols}, \emph{orthogonal}
(the effects are estimated independently, each with minimum variance).

These are the threads Chapter~\ref{ch:design} takes up in earnest---factorial and
response-surface designs, space-filling designs, and the adaptive designs that
choose each run in light of the last. The statistics of this chapter is what
makes their value precise: a design is good exactly insofar as it estimates the
effects we care about with small variance, free of bias, using the fewest runs.

\begin{figure}[t]
\centering
\includegraphics[width=0.62\linewidth]{foundations_doe.png}
\caption[Why design: OFAT versus a factorial]{A response with a genuine
interaction (its optimum lies along a diagonal ridge). One-factor-at-a-time
(red) optimises $x_1$, fixes it, then optimises $x_2$, and halts at a point that
is far from the true optimum (star)---it cannot perceive that the factors must
move \emph{together}. A $2^2$ factorial (green) samples the corners and the
centre, and the interaction is immediately visible in the responses. This is the
foundational argument for designing experiments rather than tuning factors one
at a time.}
\label{fig:stats-doe}
\end{figure}

%==============================================================================
\section{A Bayesian coda}
\label{sec:stats-bayes}

The methods above are \emph{frequentist}: a parameter is a fixed unknown and
probability describes the long-run behaviour of procedures. There is a second,
equally respectable view, and Chapter~\ref{ch:design} relies on it. The
\emph{Bayesian} treats the unknown parameter as itself uncertain and describes
that uncertainty with a probability distribution. One starts with a \emph{prior}
$p(\theta)$ encoding what is known before the data, observes data with
\emph{likelihood} $p(\mathcal D\mid\theta)$, and updates to the \emph{posterior}
by Bayes' theorem,
%
\begin{equation}
p(\theta\mid\mathcal D)\;\propto\; p(\mathcal D\mid\theta)\,p(\theta).
\label{eq:bayes}
\end{equation}
%
The posterior is the prior reweighted by how well each parameter value explains
the data. This is precisely the engine of two later constructions: the Gaussian
process surrogate of Bayesian optimisation (Section~\ref{sec:bo}), which is a
prior over \emph{functions} updated by the runs observed so far into a posterior
mean and a calibrated error bar; and the inverse parameter estimation of the
package's uncertainty-quantification layer. The two philosophies are
complementary tools, and a mature practitioner reaches for whichever answers the
question at hand.

\begin{keybox}{The thread into the next chapters}
Every statistical term in Chapters~\ref{ch:design} and~\ref{ch:ferment} is one
of the above. \emph{Least squares and ANOVA} (Sections~\ref{sec:stats-ols},
\ref{sec:stats-anova}) drive factorial response surfaces; \emph{generalized
linear models} (Section~\ref{sec:stats-glm}) analyse yes/no and count responses,
as in the strain screen; \emph{orthogonality and collinearity} explain why
designs are coded and why PCA/PLS exist; \emph{regularization}
(Section~\ref{sec:stats-reg}) makes the wide, $p>n$ screening problem tractable
and is the lens that unifies the linear models with the tree ensembles;
\emph{cross-validation} (Section~\ref{sec:stats-cv}) is the $Q^2$ and
$R^2_{\mathrm{CV}}$ that keep models honest; the \emph{Bayesian update}
(Section~\ref{sec:stats-bayes}) is Bayesian optimisation; the
\emph{variability/uncertainty} split (Section~\ref{sec:stats-varcomp}) is the
organising idea of the stochastic fermentation; and the \emph{principles of
design} (Section~\ref{sec:stats-doe}) are the reason
Chapter~\ref{ch:design} exists at all. A reader who has these in hand is ready
for the designs.
\end{keybox}
